% ----------------------------------------------------------------------
\section{Related surveys and positioning}
\label{sec:related}

Our searches identified \nSecondary{} secondary studies that touch on learning-based congestion control. \Cref{tab:surveys} compares the most relevant of them with this survey, using the scope each one states for itself. Broad surveys of machine learning (ML) or deep RL for networking~\cite{boutaba2018comprehensive,wang2018machine,luong2019applications,tang2021survey} treat congestion control as one application among many. Magazine-style overviews~\cite{zhang2019machine,zhang2020machine,wei2021congestion} discuss a handful of prominent designs and the challenges of standardizing learned controllers. Jiang et al.~\cite{jiang2021when} give the most detailed pre-2021 summary and add an experimental comparison of learning-based schemes. Huang et al.~\cite{huang2021machine} take a broad view of congestion control and avoidance that also covers routing and bandwidth reservation. Later reviews narrow the scope, to wired networks~\cite{lakshmi2022methods}, multipath TCP~\cite{maliha2023survey}, active queue management~\cite{toopchinezhad2025machine}, or fairness~\cite{haque2025fairness}. The two most recent reviews are RL-specific~\cite{li2025survey,abedini2027dynamic}. Abedini et al.\ compare RL algorithms, state/action/reward definitions, simulation platforms and reported outcomes, which is closest to our \RQ{1}. Finally, Vielhaus et al.~\cite{vielhaus2026evaluating} survey \emph{how} congestion-control algorithms in general are evaluated, which complements our \RQ{2}.

This survey differs from them as follows. (i) It is \emph{systematic}: the search, deduplication, screening and extraction are scripted and released with every decision, so the corpus can be audited and updated. (ii) It codes the \emph{evaluation level} (simulation, emulation or real network) of every study, along with generalization testing and code availability, so that claims can be weighed by their evidence. (iii) It covers all settings (Internet, wireless, cellular, satellite, IoT and data centers) through 2026. (iv) It analyzes deployment barriers, including the recent use of LLMs to design controllers, and it adds a reproducible case study on training-distribution shift.

\begin{table}[t]
\centering
\caption{Comparison with prior surveys that cover learning-based congestion control, as described by each survey's stated scope. \cmark: covered; $\circ$: partially or as one topic among many; --: not covered.}
\label{tab:surveys}
\small
\setlength{\tabcolsep}{4pt}
\begin{tabularx}{\linewidth}{@{}L{3.6cm}cYcccc@{}}
\toprule
Survey & Year & Focus & RL-specific & Wireless/DC & Eval.\ level & LLMs \\
\midrule
Luong et al.~\cite{luong2019applications} & 2019 & Deep RL in communications and networking & \cmark & $\circ$ & -- & -- \\
Zhang et al.~\cite{zhang2019machine} & 2019 & ML for Internet CC; standardization challenges & $\circ$ & -- & -- & -- \\
Zhang \& Mao~\cite{zhang2020machine} & 2020 & Selected review of ML for end-to-end CC & $\circ$ & $\circ$ & -- & -- \\
Jiang et al.~\cite{jiang2021when} & 2021 & Learning-based CC survey with experimental comparison & $\circ$ & $\circ$ & $\circ$ & -- \\
Wei et al.~\cite{wei2021congestion} & 2021 & Prominent ML-based CC redesigns compared & $\circ$ & -- & -- & -- \\
Huang et al.~\cite{huang2021machine} & 2021 & ML for broad-sense congestion control and avoidance & $\circ$ & $\circ$ & -- & -- \\
Tang et al.~\cite{tang2021survey} & 2021 & ML for end-to-end communication toward 6G & $\circ$ & $\circ$ & -- & -- \\
Lakshmi \& Arrama~\cite{lakshmi2022methods} & 2022 & RL-based CC in wired networks & \cmark & -- & -- & -- \\
Maliha et al.~\cite{maliha2023survey} & 2023 & MPTCP CC and scheduling, ML vs.\ classic & $\circ$ & $\circ$ & -- & -- \\
Toopchinezhad \& Ahmadi~\cite{toopchinezhad2025machine} & 2025 & ML for active queue management & $\circ$ & $\circ$ & -- & -- \\
Vielhaus et al.~\cite{vielhaus2026evaluating} & 2026 & How CC algorithms are evaluated (not ML-specific) & -- & \cmark & \cmark & -- \\
Li et al.~\cite{li2025survey} & 2025 & Deep RL for CC & \cmark & -- & -- & -- \\
Abedini et al.~\cite{abedini2027dynamic} & 2027 & RL-based CC: algorithms, state/action/reward, platforms & \cmark & $\circ$ & $\circ$ & -- \\
\midrule
\textbf{This survey} & 2026 & Systematic review of RL/IL CC across all settings & \cmark & \cmark & \cmark & \cmark \\
\bottomrule
\end{tabularx}
\end{table}
